\section{总结与延伸}

\subsection{核心结论}
\begin{importantbox}{本讲可执行结论}
\begin{enumerate}
    \item 先建立真实交互评估，再谈模型优化；
    \item 推理时搜索是长链路任务鲁棒性的核心；
    \item Verifier/Judge 是连接搜索、训练和数据过滤的关键组件；
    \item 数据闭环决定能力能否持续增长；
    \item 失败恢复机制与人机协同必须在系统设计早期纳入。
\end{enumerate}
\end{importantbox}

\subsection{总结表}
\begin{table}[H]
\centering
\small
\begin{tabular}{p{2.6cm}p{3.9cm}p{4.3cm}}
\toprule
\textbf{主题} & \textbf{讲者重点} & \textbf{工程动作} \\
\midrule
评估 & VisualWebArena 衡量真实交互能力 & 建立终局 + 轨迹 + 错误类型联合指标 \\
搜索 & 单轨执行不稳，需可回溯决策 & 上线分层搜索与预算控制 \\
验证器 & Judge 更擅长 ``评'' 而非 ``做'' & 把 verifier 前移到推理中间步骤 \\
数据扩展 & 在线采样 + 过滤形成闭环 & 构建失败样本回流与难度重加权 \\
鲁棒性 & 失败恢复与人工协同不可省略 & 增加回滚、重规划、澄清触发机制 \\
\bottomrule
\end{tabular}
\caption{Multimodal Autonomous AI Agents 讲座的工程落地摘要}
\end{table}

\subsection{本章小结}
这节课给出的不是单一算法，而是一套 Agent 工程方法论：以真实评估为锚点，用搜索和验证器提升推理鲁棒性，用数据闭环维持长期迭代速度。

\subsection{Further Reading}
\begin{itemize}
    \item VisualWebArena / WebArena 相关论文与基准代码；
    \item LLM-as-a-Judge 与 trajectory-level verifier 研究；
    \item Search-Augmented Inference 在 Agent 任务上的方法比较；
    \item 在线交互数据收集与高置信度轨迹过滤实践；
    \item 多智能体规划与并行执行框架（用于高复杂度长任务）。
\end{itemize}

\end{document}
